\documentclass[a4]{article}
\usepackage{pack}
% Pacchetti specifici per gli appunti
\usepackage{todonotes}
\usepackage{parskip} % Evita che ci sia un rientro nei nuovi paragrafi

\geometry{
	%twocolumn,
	a4paper,
	total={170mm,257mm},
	left=30mm,
	right=30mm,
}
\title{Caratterizzazione di dinamiche markoviane in regime quantistico \\[1ex]
	   \large{Note e riflessioni}}
\author{Luca Arnaboldi \\ 
	
	\includegraphics[width=0.2\textwidth]{img/unipi-noscritta.png}}
\date{17 settembre 2020}

\begin{document}
\maketitle
\tableofcontents
\clearpage
Il setting è quello degli \emph{open quantum systems}. In particolare vogliamo studiare
la markovianità dell'evoluzione temporale, spesso chiamata \emph{divisibilità}, anche se
le due definizioni non coincidono. Lo scopo della tesi è indagare le differenze tra questi
due concetti. Il lavoro si basa sul contenuto degli articoli \emph{Complete Positive
Divisibility does not mean Markovianity} di Modi~et~al \cite{CPdoesnotimply}, e \emph{Quantum non-Markovian processes break conditional past-future independence} di Budini \cite{budini2018quantum}.

\section{Notazione}
Descriveremo un sistema \(S\), a contatto con un ambiente \(E\); i rispettivi spazi di
Hilbert sono \(\Hilb_S\) e \(\Hilb_E\), supposti tutti di dimensione finita. Lo spazio
su cui si lavora è dunque 
\(\Hilb_{SE} = \Hilb_{SE} \otimes \Hilb_{SE}\), e la più generica Hamiltoniana si
scriverà come
\[\hat{H}_{SE} = \hat{H}_S \otimes  \Id_E + \Id_S \otimes  \hat{H}_E + \hat{V}_{SE}.\]

L'insieme degli operatori lineari di \(\Hilb\) è \(\LinSet(\Hilb)\), mentre quello delle
matrici densità è \(\sigma(\Hilb)\); ovviamente \(\LinSet(\Hilb) \subset \sigma(\Hilb)\).
Indicheremo con \(\hat{\rho}_S\) e \(\hat{\tau}_E\) due generiche matrici densità di \(S\)
e \(E\) rispettivamente; la matrice di densità completa è invece indicata da \(\hat{\rho}_{SE}\).
Ovviamente tutte dipendono dal tempo.


Indicheremo con \(\Phi_{t'\to t}(\cdot)\) la mappa che evolve che manda \(\hat{\rho}_S(t')\)
in \(\hat{\rho}_S(t)\). Se omettiamo \(t'\) significa che sottointendiamo che parte da \(t'=0\).
Per brevità, a volte indicheremo \(\Phi_{t'\to t}\) con \(\Phi_{t:t'}\) (attenzione il verso è scambiato!).
Infine se la mappa \(\Phi\) è accessibile sperimentalmente verrà indicata con
\(\Lambda\).

\section{Ripasso}
Qui metto un po' di cose ``scolastiche'', cioè prese dai corsi di \texttt{MQ} o \texttt{QINFO}
che però ci tengo a ripassare per la tesi.
\subsection{Misure \& POVM}
Ho preso tutto dal Nielsen \cite{nielsen2010quantum}.

Il postulato di misura della meccanica quantistica (\emph{misure proiettive}) è particolarmente
restrittivo. Ad esempio non contempla casi in cui la misura è irripetibile: se ripeto una
seconda volta una misura proiettiva non perturbo nuovamete il sistema ed ottengo lo stesso
risultato. Chiaramente però questo tipo di misura non può ovviamente descrivere il caso in
cui, ad esempio, ho misurato un fotone distuggendolo.
\begin{post}[di misura]
Le misure in meccanica quantistica sono descritte da un insieme di operatori
\(\{\hat{M}_m\}\), tali che rispettano l'equazione di completezza
\begin{equation} \label{eq:compmis}
\sum_m \hat{M}^\dag_m \hat{M}_m = \Id
\end{equation}
e l'indice \(m\) si riferisce al risultato della misura.

Se il sistema si trova nello stato \(\ket{\psi}\) prima della misura, si troverà successivamente
nello stato
\[ \frac{\hat{M}_m \ket{\psi}}{\sqrt{\Braket{\psi|\hat{M}^\dag_m \hat{M}_m|\psi}}}\]
con probabilità \(p(m)= \braket{\psi|\hat{M}^\dag_m \hat{M}_m|\psi}\).
\end{post}
L'equazione \eqref{eq:compmis} discende dalla richiesta che le probabilità sommino a 1:
\(\sum_m~p(m)~=~1\).

I POVM sono utili quando non siamo interessati allo stato del sistema dopo la misura.
\begin{defn}[POVM]
Un \emph{POVM} è un insieme di operatori \emph{positivi} \(\{\hat{E}_m\}\) tali che
\[\sum_m \hat{E}_m  = \Id.\]
\end{defn}
È una buona definizione perché scegliendo \(\hat{M}_m = \sqrt{\hat{E}_m}\) otteniamo una
misura valida per il postulato di misura. I POVM forniscono infomazioni \emph{solo sulla
probabilità di misura} e non sullo stato risultante siccome comunque scegliendo
\[\hat{M}_m = \hat{U}\sqrt{\hat{E}_m},\]
dove \(\hat{U}\) è una qualsiasi unitaria, si ottiene una misura con la stessa statistica
del POVM, ma che ovviamente non ha lo stesso stato finale. Questa ambiguità
rende i POVM inutilizzabili nei casi in cui è necessario conoscere lo stato.

Una misura proiettiva si ottiene scegliendo come come operatori di misura
i poiettori. Analogamente la descrizione POVM di una misura poiettiva è
\[\hat{E}_m = \ket{m}\!\bra{m}.\]

Per concludere, la probabilità di ottenere il risultato \(m\) da una misura su
una matrice di densità è
\[p(m) = \Tr[\hat{E}_m \hat{\rho}]\]



\subsection{Open Quantum Systems: cose generali}
Per le note proprietà delle matrici di densità vale 
\[\hat{\rho}_S(t) = \Tr_E\left(\hat{\rho}_{SE}(t)\right).\]
Ovviamente il sistema congiunto \(S-E\) evolve come un sistema chiuso, quindi attraverso un
operatore unitario \(\hat{U}_{SE}(t, t')\):
\[\hat{\rho}_{SE}(t) = \hat{U}_{SE}(t, t') \hat{\rho}_{SE}(t') \hat{U}^{\dag}_{SE}(t, t'),\]
ma lo stesso non vale per il sistema \(S\) da solo. Infatti lui non è chiuso, dato che interagisce
con l'ambiente; \(S\) è appunto detto \emph{open quantum system}.
\begin{figure}[h!]
	\centering
	\resizebox{0.5\textwidth}{!}{\begin{tikzpicture}
\draw[rounded corners, dashed] (0.75,0.75) rectangle (4.75,4.77);
\filldraw[rounded corners, fill=pastelgreen] (1,1) rectangle (4.5,2);
\draw[double arrow=3pt colored by black and pastelred]
    (2.5,2.05) arc(210:150:.25) arc(-30:30:.25) arc(210:150:.25) arc(-30:30:.25)-- ++(90:.5);
\draw[double arrow=3pt colored by black and pastelred]
    (3,3.50) arc(150:210:.25) arc(30:-30:.25) arc(150:210:.25) arc(30:-30:.25)-- ++(270:.5);
\filldraw[rounded corners, fill=pastelblue] (2.25,3.52) rectangle (3.25,4.52);

\node[white] at (2.75,4.02) {\(S\)};
\node[white] at (2.75,1.5) {Ambiente \(E\)};
%    \filldraw[rounded corners, fill=pastelyellow] (0,0) rectangle (6,6);
%    \filldraw[fill=pastelgreen] (3,3) circle (1.3cm);
%    %E
%    \draw[double arrow=3pt colored by black and pastelred]
%    (4.4,3) arc(120:60:.25) arc(-120:-60:.25)arc(120:60:.25) arc(-120:-60:.25) -- ++(0:.5);
%    %N
%    \draw[double arrow=3pt colored by black and pastelred]
%    (3,4.4) arc(210:150:.25) arc(-30:30:.25) arc(210:150:.25) arc(-30:30:.25)-- ++(90:.5);
%    %W
%    \draw[double arrow=3pt colored by black and pastelred]
%    (1.6,3) arc(60:120:.25) arc(-60:-120:.25)arc(60:120:.25) arc(-60:-120:.25) -- ++(180:.5);
%    %S
%    \draw[double arrow=3pt colored by black and pastelred]
%    (3,1.6) arc(150:210:.25) arc(30:-30:.25) arc(150:210:.25) arc(30:-30:.25)-- ++(270:.5);
    
\end{tikzpicture}}
	\caption{open  quantum system.}
\end{figure}
\begin{defn}[Completamente Positivo] Una mappa \(\Phi\colon \sigma(\Hilb) \to \sigma(\Hilb) \)
si dice \emph{completamente positiva} se dato un qualsiasi sistema \(A\), la mappa
\[\Phi \otimes \Id_A: \sigma(\Hilb \otimes \Hilb_A) \to \sigma(\Hilb \otimes \Hilb_A) \]
è positiva.
\end{defn}

\(\Phi_{t'\to t}(\cdot)\) è lineare, completamente positiva e conserva la traccia,
più brevemente  diremo \emph{LCPT}. CP è vero nel caso in cui lo stato iniziale di
S-E è fattorizzabile (vedi tipo \cite{breuer2007theory} per i dettagli).

\begin{defn}[Divisibilità completamente positiva] Un sistema è \emph{CP-divisibile} se la sua
mappa \(\Phi\) è tale che
\begin{equation} \label{eq:CPdiv}
\Phi_{t'\to t} = \Phi_{s\to t} \circ \Phi_{t'\to s} \quad \forall t, t', s: t' \le s \le t.
\end{equation}
\end{defn}

\subsection{Isomorfismo di Choi–Jamiołkowski}
Consideriamo un canale quantistico \(\Phi: \LinSet(\Hilb_\mathcal{S}) \to
\LinSet(\Hilb_\mathcal{S'})\). Di solito (tipo a \texttt{QINFO2}) si considera un canale di
un sistema in sè stesso, ma dato che nell'articolo lo definiamo così, allora facciamo tra
due spazi diversi.
\begin{defn}[Stato di Choi–Jamiołkowski di \(\Phi\)]
Sia \(\mathcal{A}\) un sistema ausiliario (\emph{ancilla}) della stessa dimensione di
\(\mathcal{S}\). Siano \( \left\{\ket{i}_\mathcal{S}\right\}_{i=1,\dots,d} \)  e
\(\left\{\ket{i}_\mathcal{A}\right\}_{i=1,\dots,d}\) basi ortonormali di \(\Hilb_\mathcal{S}\) e
\(\Hilb_\mathcal{A}\) rispettivamente. Si dice \emph{stato di Choi–Jamiołkowski di \(\Phi\)}
\[\ChoiState{\Phi} := \left( \Phi \otimes \Id_\mathcal{A}\right) 
  \left(\sum_{i=1}^d \frac{\ket{i}_\mathcal{S} \ket{i}_\mathcal{A} \bra{i}_\mathcal{S}
  \bra{i}_\mathcal{A}}{\sqrt{d}}\right) \in \sigma(\Hilb_\mathcal{S'} \otimes \Hilb_\mathcal{A}).\]
\end{defn}
Facciamo ora qualche osservazione. La definizione \emph{non è indipendente dalle basi} scelte:
per scelte diverse degli stai \(\ket{i}\) si ottengono stati diversi. In secondo luogo a priori
non è per nulla chiaro che \(\ChoiState{\Phi}\) sia effettivamente una buona matrice 
di densità. Sappiamo solo che è un operatore positivo dato che \(\Phi\) è CP. Si dimostra però
(vedi sul quaderno di \texttt{QINFO2}) che
\[\Phi\text{ LCP} \quad \implies \ChoiState{\Phi} \in \sigma(\Hilb_\mathcal{S'}
\otimes \Hilb_\mathcal{A}).\]
Inoltre le peculiarità non finiscono qui, dato che vale il seguente teorema.
\begin{thm}[di Choi sulle mappe CP] \label{thm:choi}
L'operatore \(\ChoiState{\Phi}\) è semi-definito positivo se e solo se \(\Phi\) è una mappa
completamente postiva.
\end{thm}
Da notare che questo fornisce un ottimo criterio per verificare la positività completa di
una mappa.

Notiamo infine che non è un isomorfismo nel senso che ci aspettiamo: non è vero che ad ogni
elemento di \(\sigma(\Hilb_\mathcal{S'}\otimes \Hilb_\mathcal{A})\) corrisponde una mappa
LCPT. Infatti vale il criterio
\[\Phi\text{ LCPT} \iff \Tr_\mathcal{S}\left[\ChoiState{\Phi}\right]=\frac{\Id_{\mathcal{A}}}{d}.\]
L'isomorfismo vale nel senso del Teorema~\ref{thm:choi}: ad ogni mappa CP corrisponde
(a meno della normalizzazione) una matrice densità.

\subsection{Master equation}
\begin{defn}[Semigruppo dinamico quantistico]
	Un canale quantistico si dice essere un semigruppo dinamico se è CP-divisibile
	ed omogeneo nel tempo, cioè
	\[\Phi_{t'\to t} \equiv \Phi_{t'- t}.\]
	Si può dunque scrivere
	\[\Phi_{t} \circ \Phi_{s} = \Phi_{t+s}.\]
\end{defn}
A lungo si è considerato questa la definizione di markovianità per un sistema quantistico
aperto, o almeno una condizione sufficiente; il corso di \texttt{QINFO 2} diceva proprio questo, ma è l'obbiettivo di questo lavoro dimostrare che non è così.

Vale il seguente teorema
\begin{thm}[GKSL\footnote{Gorini–Kossakowski–Sudarshan–Lindblad, almeno una volta
		                  andava scritto.}] \label{thm:GKSL}
	Un sistema quantistico forma un semi-gruppo dinamico se e solo se vale
	\[\dot{\rho}_S(t)=\mathcal{L}\left(\rho_S(t)\right),\]
	con \(\mathcal{L}\) detto \emph{lindbladiano} della forma
	\[\mathcal{L}\left[\hat{\theta}\right] 
	  = -i\left[\hat{H},\hat{\theta}\right]+
	    \sum_k \left(\left[\hat{L}_k\hat{\theta},\hat{L}^\dag_k\right]+
	                 \left[\hat{L}_k,\hat{\theta}\hat{L}^\dag_k\right] \right).\]
	\(\hat{H}\) è hermitiano, mentre \(\left\{\hat{L}_k\right\}\) è detto insieme
	degli operatori di Lindeblad.
\end{thm}
Un processo però può essere CP divisibile, ma non essere omogeneo.
Nel riferimento~\cite{breuer2016colloquium} si dice che un processo continua a essere
CP divisibile se il limbladiano è della forma
\[\mathcal{L}(t)\left[\hat{\theta}\right] 
= -i\left[\hat{H},\hat{\theta}\right]+
\sum_k \gamma_i(t) \left(\left[\hat{L}_k\hat{\theta},\hat{L}^\dag_k\right]+
\left[\hat{L}_k,\hat{\theta}\hat{L}^\dag_k\right] \right)
\quad \text{con} \quad \gamma_i(t) \ge 0 \quad \forall i, \forall t.\]
\begin{tcolorbox}[colback=yellow!5!white,colframe=yellow!85!black,title=Nota]
	Non sono riuscito a capire se questa è una condizione sufficiente. Nell'articolo
	\cite{breuer2016colloquium} in particolare non si capisce se si possono avere anche
	dipendenze temporali su gli \(\left\{\hat{L}_k\right\}\) e \(\hat{H}\). Secondo me
	su quest'ultimo sì, mentre su gli altri no, ma non sono sicuro.
	
	Quella che ho scritto io è sicuramente una condizione sufficiente, è confermato anche
	da \cite{budini2018quantum} nel primo esempio.
\end{tcolorbox}

\subsection{Markovianità classica}
Consideriamo solo processi discreti, tanto le misure di sistemi quantistici le possiamo 
fare solo a  tempi distinti. Baso le definizione sul Norris~\cite{norris1998markov}, che
può essere consultato  per saperne di più.

Indicheremo con \(\{X_n\}\) le variabili casuali di un processo. Intuitivamente un processo
è markoviano  quando lo step successivo dipende solo dalla storia precedente, e non da
tutta la storia passata 
\begin{defn}[Processo Markoviano] \label{defn:classicalmarkov}
Diremo che \(\{X_n\}\) è  un processo markoviano con distribuzione iniziale \(\lambda\) se
\[\mathbb{P}\left(X_{n+1} = a_{n+1} | X_{n} = a_{n}, \dots, X_0 = a_0\right) =
  \mathbb{P}\left(X_{n+1} = a_{n+1} | X_{n} = a_{n}\right)=P^{(n)}_{a_na_{n+1}} \quad \forall n\in \mathbb{N}.\] 
Diremo per brevità  che il processo è \(\text{Markov}(\lambda,P)\).
\end{defn}
Si è sottolineata la dipendenza di \(P\) da \(n\), nel senso che un processo resta markoviano
anche se le probabilità di transizione cambiano nel tempo. Se invece non succede allora il
processo si dice anche \emph{omogeneo}.

Esponiamo ora un importante risultato  che permette di dare una definizione equivalente di processo markoviano.
\begin{thm} \label{thm:markovalt}
Un processo è \(\text{Markov}(\lambda,P)\) se e solo se 
\[\mathbb{P}\left(X_0 = a_0, X_1 = a_1, \dots, X_n = a_n\right) =
  \lambda_0 P^{(0)}_{a_0a_1} P^{(1)}_{a_1a_2} \cdots P^{(n-1)}_{a_{n-1}a_n}
  \quad \forall n\in \mathbb{N}.\]
\end{thm}




\section{Inquadriamo il problema}
La definizione classica di processo markoviano è ben nota (vedi tipo \cite{breuer2007theory, markovcondition}).
Il problema se consideriamo processi quantistici è che, nonostante sia matematicamente
ben definita, \emph{operativamente} è inutilizzabile: non posso infatti verificare
la markovianità (che è una condizione ad  ogni tempo) senza perturbare il sistema
quando misuro. In ultima analisi quindi  la definizione classica di markovianità
\textbf{non} è utilizzabile nel linguaggio quantistico.
La gente ha fatto diversi tentativi per risolvere il problema, ma quasi tutti seguono
il pattern: prendo una condizione necessaria, \emph{ma non sufficiente} alla markovianità
classica, e la estendo facendola diventare una condizione quantistica. La CP-divisibilità
è uno di questi tentativi: ho preso l'equazione di Chapman-Kolmogorov (che è in un certo
senso la proprietà più importante dei processi markoviani), e l'ho trasformata in una
condizione  quantistica.
È stata  trovata una buona definizione di markovianità quantistica, che coincide con
quella classica per il limite giusto, è tutto esposto nel \cite{markovcondition}.
Il problema è dunque capire come si inquadra in questo  contesto la completa positività.

I problemi però non finiscono qui: l'attuale definizione di CP-positività non si presta alla
verifica in laboratorio  (nel senso  che manca una definizione precisa su che tomografia
quantistica effettuare per verificare la completa positività), dovremo dunque occuparci
di formularla in una definizione operativa.
Infine vedremo come metterla in relazione con la markovianità quantistica.

\section{Markovianità quantistica}
Espongo in questa sezione due definizioni della markovianità quantistica differenti,
proposte da diversi autori.
\subsection{Causal break \textit{di Modi} \cite{markovcondition}}
La definizione è riportata nell'articolo \cite{markovcondition} e non la trascrivo qui
perché non ne ho voglia. Essa discende dalla definizione~\ref{defn:classicalmarkov},
che viene recuperata nel limite classico.
\subsection{Past-future correlations \textit{di Budini} \cite{budini2018quantum}}
La definizione discende dalla definizione classica equivalente data dal
Teorema~\ref{thm:markovalt}. Infatti usando quest'ultimo si può dire
\begin{align*}
	\mathbb{P}(X_{n-1}=a_{n-1}, X_{n+1}=a_{n+1}|X_{n}=a_{n})&=
    \frac{\mathbb{P}(X_{n-1}=a_{n-1}, X_{n}=a_{n}, X_{n+1}=a_{n+1})}{\mathbb{P}(X_{n}=a_{n})}=\\
  &=\frac{\mathbb{P}(X_{n-1}=a_{n-1}, X_{n}=a_{n}) P^{(n)}_{a_n a_{n+1}}}
  {\mathbb{P}(X_{n}=a_{n})}=\\
  &=\frac{\mathbb{P}(X_{n-1}=a_{n-1}, X_{n}=a_{n}) \mathbb{P}(X_{n+1}=a_{n+1}|X_{n}=a_{n})}
  {\mathbb{P}(X_{n}=a_{n})}=\\
  &=\mathbb{P}(X_{n-1}=a_{n-1}|X_{n}=a_{n}) \mathbb{P}(X_{n+1}=a_{n+1}|X_{n}=a_{n})
\end{align*}
Possiamo dunque usare  la quantità\footnote{Ho droppato un bel po' di notazione perché
mi sono rotto di scrivere tutta quella roba, tanto è chiaro l'intento. Inoltre, formalmente
dovrei anche dire che la il processo è \emph{omogeneo} (cioè \(P\) non dipende da
\(n\)), però anche nell'articolo  se ne fregano. Infine la definizione in linea teorica
dipende dallo step in mezzo, ma anche qui l'articolo glissa, e io mi comporterò in maniera
simile.} 
\[C_{pf} = \sum_{x,z}\left(P(x,z|y)-P(x|y)P(z|y)\right)\]
detta \emph{past future correlation}, per discriminare i processi markoviani da
quelli non-markoviani: infatti come mostrato sopra per un processo markoviano si ha
\(C_{pf} = 0\). La cosa si estende facilmente anche a casi dove al
posto che fissare un singolo step se ne  fissa più che uno, diciamo \(k\)
\[C^{(k)}_{pf} = \sum_{x,z}\left(P(x,z|\bm{y})-P(x|\bm{y})P(z|\bm{y})\right),\]
dove però ora \(\bm{y}\) è un vettore di dimensione \(k\).
Per un processo markoviano \(C^{(k)}_{pf} \quad \forall k\).
\begin{tcolorbox}[colback=yellow!5!white,colframe=yellow!85!black,title=Nota]
	Nell'articolo sembra essere detto implicitamente che se \(C^{(k)}_{pf}\)
	è nullo per ogni \(k\) allora il  processo è markoviano. Purtroppo non c'è traccia
	di dimostrazione e a me sembra proprio una  di quelle cose che vorresti che siano vere
	ma in realtà sono false. Ad esempio questo \url{https://math.stackexchange.com/questions/3147222/if-all-powers-of-two-random-variables-are-uncorrelated-are-they-independent}
\end{tcolorbox}


\section{Classi di divisibilità}
Esponiamo innanzi tutto due definizioni non equivalenti della \emph{CP-divisibilità}.

\subsection{iCP divisibilità}
Siamo in un contesto sperimentale dove possiamo preparare un sistema in un qualsiasi
stato iniziale al tempo \(r=0\), e poi effettuare misure a qualsiasi tempo \(s>0\).
Con la tomografia  quantistica  possiamo costruire  l'insieme
\[  \lambda_0 = \{\Lambda_{s\colon0}\}, \]
delle mappe che descrivono la dinamica. Supponiamo che \textbf{tutte le mappe di
\(\boldmath{\lambda_0}\) siano invertibili.}
\begin{defn}[iCP divisibilità]
Un processo è \emph{CP-divisibile per inversione} se, dati \(s \le  t\), la mappa
\[\Phi_{t\colon s} = \Lambda_{t\colon0} \circ \Lambda^{-1}_{s\colon0}\]
è completamente positiva.
\end{defn}
\begin{figure}
	\centering
	\resizebox{0.4\textwidth}{!}{
\begin{tikzpicture}
\filldraw[rounded corners, fill=pastelblue] (-1,0) -- (1,0) -- (0,-1) -- cycle;
\filldraw[rounded corners, fill=pastelgreen] (-0.5,3) -- (-0.5,1) -- (-1.5,2) -- cycle;
\draw [rounded corners,line width=1mm, black] (0,0) -- (0,1) -- (1.5,1);
\draw [line width=1mm, black] (-0.5,2) -- (1.5,2);
\filldraw[rounded corners, fill=pastelred] (1.5,0.5) rectangle (3,2.5);
\draw [line width=1mm, black] (3,2) -- (3.5,2);
\draw [rounded corners,line width=1mm, black] (3.2,1.7) -- (3.8,2.3);
\draw [line width=1mm, black] (3.0,1) -- (3.3,1);

\node[white] at (-0.9,2) {\LARGE \(\hat{\eta}_r\)};
\node[white] at (0,-0.4) {\LARGE \(\hat{\rho}_r\)};
\node at (3.6,1) {\LARGE \(\hat{\rho}_{s}\)};

\node[white] at (2.3,1.5) {\LARGE \(\hat{U}_{s\colon r}\)};
\end{tikzpicture}
}
	\caption{schema rappresentante un processo iCP.}
\end{figure}
Ovviamente se un processo è iCP divisibile allora ogni mappa \(\Phi_{t\colon s}\) è
ben definita a partire da \(\lambda_0\).

Questa è talvolta considerata la definizione di CP divisibilità, vedi ad esempio
\cite{breuer2016colloquium}. Questo lo si fa perché ogni canale quantistico che
è invertibile è scrivibile i una forma analoga a quella del Teorema~\ref{thm:GKSL},
dove però gli \(\left\{\hat{L}_k\right\}\) e \(\hat{H}\) possono dipendere dal tempo.
Ovviamente non è più vero che questo è un semi-gruppo dinamico in caso di dipendenza
temporale.

\subsection{oCP divisibilità}
Supponiamo invece che di poter manipolare il sistema a qualsiasi tempo \(s\). In
particolare possiamo interrompere l'evoluzione ad un istante \(s_-\), effettuare una
misura, e sostituire lo stato con una qualsiasi altra matrice di densità in modo che
l'evoluzione continui dall'istante \(s_+\). In questa situazione è  possibile  ricostruire
sperimentalmente tutte le mappe dell'insieme
\[\lambda = \{\Lambda_{t\colon s}\}  \quad \forall s,t\colon s \le t.\]  
\begin{defn}[oCP  divisibilità]
Un processo è \emph{CP-divisibile operativamente} se dati \(r \le s \le  t\), vale la
relazione 
\begin{equation} \label{eq:OCPdiv}
\Lambda_{t\colon r} = \Lambda_{t\colon s} \circ \Lambda_{s\colon r},
\end{equation}
dove le mappe appartengono tutte all'insieme \(\lambda\).
\end{defn}
\begin{figure}
	\centering
	\resizebox{0.9\textwidth}{!}{
\begin{tikzpicture}
\filldraw[rounded corners, fill=pastelblue] (-1,0) -- (1,0) -- (0,-1) -- cycle;
\filldraw[rounded corners, fill=pastelgreen] (-0.5,3) -- (-0.5,1) -- (-1.5,2) -- cycle;
\draw [rounded corners,line width=1mm, black] (0,0) -- (0,1) -- (1.5,1);
\draw [line width=1mm, black] (-0.5,2) -- (1.5,2);
\filldraw[rounded corners, fill=pastelred] (1.5,0.5) rectangle (3,2.5);
\draw [rounded corners,line width=1mm, black] (3,1) -- (3.5,1) -- (3.5,0.);
\draw [rounded corners,line width=1mm, black] (3.2,-0.3) -- (3.9,0.4);
\filldraw[rounded corners, fill=pastelblue] (4.5,0) -- (6.5,0) -- (5.5,-1) -- cycle;
\draw [rounded corners,line width=1mm, black] (5.5,0) -- (5.5,1) -- (7,1);
\draw [line width=1mm, black] (3,2) -- (7,2);
\filldraw[rounded corners, fill=pastelred] (7,0.5) rectangle (8.5,2.5);
\draw [line width=1mm, black] (8.5,2) -- (9,2);
\draw [rounded corners,line width=1mm, black] (8.7,1.7) -- (9.3,2.3);
\draw [line width=1mm, black] (8.5,1) -- (9.3,1);

\node[white] at (-0.9,2) {\LARGE \(\hat{\eta}_r\)};
\node[white] at (0,-0.4) {\LARGE \(\hat{\rho}_r\)};
\node at (3.55,-0.7) {\LARGE \(\hat{\rho}_{s_-}\)};
\node[white] at (5.5,-0.4) {\LARGE \(\hat{\rho}_{s_+}\)};
\node at (9.6,1) {\LARGE \(\hat{\rho}_{t}\)};

\node[white] at (2.3,1.5) {\LARGE \(\hat{U}_{s\colon r}\)};
\node[white] at (7.75,1.5) {\LARGE \(\hat{U}_{t\colon s}\)};

\end{tikzpicture}
}
	\caption{schema rappresentante un processo oCP.}
\end{figure}
Le equazioni \eqref{eq:CPdiv} e \eqref{eq:OCPdiv} sono formalmente identiche, con la
differenza che nella seconda è chiaro come verificarla operativamente.

\paragraph{Conditional non-signalling}
La definizione di oCP divisibilità è ancora ambigua, dato che la mappa \(\Lambda_{t\colon s}\)
potrebbe a priori dipendere da \(\hat{\rho}^{(r)}\), ovvero lo stato iniziale del sistema.
In effetti questo non può succedere, come mostriamo nel Lemma~\ref{lem:condnsig}.
\begin{lem}  \label{lem:condnsig}
	Siano \(r<s<t\) tre istanti temporali. L'\emph{oCP divisibilità} implica che non ci
	sia \emph{signalling} tra \(r\) e \(t\), \emph{condizionato} al fatto che all'istante
	\(s_-\) lo stato del sistema \(S\) è scartato e preparato in uno nuovo all'istante \(s_+\).
	
	\proof Basta dire che se così non fosse la oCP divisibilità non sarebbe una buona
	definizione. 
	
	Più formalmente si può osservare che se non valesse il \emph{conditional non-signalling}
	allora la mappa \(\Lambda_{t\colon s}\) non potrebbe essere CP. 
	Supponiamo per assurdo che la dinamica sia conditionally signalling. Allora esiste
	uno stato \(\hat{\rho}_s\) nel quale possiamo preparare il sistema al tempo \(s\)
	tale che lo stato finale al tempo \(t\) dipende dalla scelta iniziale dello stato
	 \(\hat{\rho}_r\). Se così fosse la mappa \(\Lambda_{t\colon s}\) non potrebbe essere
	 CP dato che per lo stesso input \(\hat{\rho}_s\) dovrebbe avere output diversi, e ciò
	 è impossibile per le mappe CP.
	 \todo{Ma non dovrebbe essere impossibile per le funzioni in genelare?}
\end{lem}
\begin{ach}
	Non vale il viceversa, ovvero esistono processi \emph{conditionally non-signalling}, ma
	che non sono oCP divisibili. Un esempio è si trova nel paragrafo \ref{cont:cnsocp}.
\end{ach}

\subsubsection{Conditional signalling formalmente}
Nell'Appendice~A dell'articolo~\cite{CPdoesnotimply} c'è una definizione formale di
\emph{non-signalling} e \emph{conditional non-signalling} in termini di POVM e preparazioni
tomograficamente complete. Non la trascrivo perché è già chiaro da là.

\subsection{Relazioni tra le classi di divisibilità}
Come specificato in precedenza l'iCP divisibilità ha senso fin tanto che le mappe in gioco
sono invertibili. Ristretti a questo caso è banale vedere che oCP~\(\implies\)~iCP, basta 
moltiplicare a desta per l'inverso di \(\Lambda_{s\colon r}\) e si ottiene l'iCP divisibilità.
Si potrebbe pensare che vale il viceversa, ma nella sezione~\ref{cont:icpocp} si mostra
un processo iCP, ma che è allo stesso tempo conditional signalling.

L'oCP divisibilità non è comunque sufficiente ad eliminare tutti gli effetti di memoria
dell'ambiente: esistono infatti processi che sono oCP divisibili, ma non markoviani.
Si veda il controesempio~\ref{cont:ocpmar} per un esempio esplicito.
L'oCP divisibilità però corrisponde a quella che può essere chiamata \emph{markovianità
i media}, cioè lo stato finale dipende solo dall'ultima preparazione a patto di mediare
su tutte le misure passate.
\begin{figure}[h!]
	\centering
	\resizebox{0.5\textwidth}{!}{\begin{tikzpicture}
	
	\fill[rounded corners,fill=pastelblue] (0,0) rectangle (6.5,4.5);
    \draw[color=black,dashed] (3,5.0) -- (3,-0.5);
    \fill[rounded corners,fill=pastelgreen] (0.5,0.25) rectangle (3,3.5);
    \fill[rounded corners,fill=pastelyellow] (1,0.35) rectangle (5.5,2.5);
    \fill[rounded corners,fill=pastelred] (1.5,0.45) rectangle (5.0,1.5);
    
    \node[white] at (3.25,1) {Markovianità};
    \node at (3.25,2) {oCP divisibilità};
    \node[white] at (1.75,3) {iCP divisibilità};
    \node[white] at (3.25,4) {Non divisibilità};
    \node at (1.625,4.8)  {\(\Lambda_{t:r}\) invertibile};
    \node at (4.875,4.8)  {\(\Lambda_{t:r}\) non invertibile};
\end{tikzpicture}}
	\caption{classi di divisibilità e markovianità.}
\end{figure}

%% Per fare davvero l'esempio dovrei spiegare tutta la roba di come trasformare il processo
%% in una statistica classica quando misuro. Probabilmente non ne ho voglia perchè è fatto
%% bene sull'articolo.

%Per essere più espliciti consideriamo come al solito un processo che inizia al tempo
%\(r\), viene misurato al tempo \(s_-\), preparato nuovamente ad \(s_+\) ed infine misurato
%defitivamente al tempo \(t\); indichiamo con \(\hat{\rho}^{(s_-)}\) e \(\hat{\rho}^{(t)}\)
%gli stati misurati al tempo \(s_-\) e \(t\) rispettivamente e con \(\hat{\rho}^{(r)}\) e
%\(\hat{\rho}^{(s_+)}\) l'ensemble statistico degli stati misurati ai tempi \(r\) e
%\(t\) rispettivamente.
Un processo è markoviano invece quando lo stato finale dipende solo dall'ultima trasformazione
\emph{per ogni} sequenza di misure precedenti.






\section{Controesempi}
\subsection[Cond. non-signalling does not imply oCP]
           {Conditionally non-signalling\(\centernot\implies\)oCP} \label{cont:cnsocp}
Consideriamo un sistema dove l'ambiente è costituito da due quibit: \(\mathcal{E}_s\) e
\(\mathcal{E}_r\). Il sistema d'interesse \(\mathcal{S}\) è anch'esso un qubit.
Nell'istante iniziale \(r\) l'ambiente si trova nello stato
\[\hat{\eta}_{\mathcal{E}_s\mathcal{E}_r} =
  \frac{1}{4}\left(\Id_{\mathcal{E}_s\mathcal{E}_r} +
                    \SigmaZ_{\mathcal{E}_s} \otimes \SigmaX_{\mathcal{E}_r}
             \right). \]
Indicheremo con \(\hat{\rho}^{(r)}\) lo stato in cui è preparato il sistema al tempo \(r\).
Siamo interessati all'evoluzione in due istanti differenti \(r<s<t\).

L'evoluìone temporale è descritta da
\begin{align*}
	\hat{U}_{\mathcal{S}\mathcal{E}_s\mathcal{E}_r}(r \to s) &= 
	  \hat{S}_{\mathcal{S}\mathcal{E}_r} \otimes \Id_{\mathcal{E}_s} \\
	\hat{U}_{\mathcal{S}\mathcal{E}_s\mathcal{E}_r}(s \to t) &=
	  \hat{U} = \frac{1}{\sqrt{3}}\left(\Id_{\mathcal{S}\mathcal{E}_s} +
	    	                            i \SigmaY_{\mathcal{S}} \otimes \Id_{\mathcal{E}_s} +
	    	                            i \SigmaY_{\mathcal{S}} \otimes \SigmaZ_{\mathcal{E}_s}
	    	                      \right) \otimes \Id_{\mathcal{E}_r},    	    	                              
\end{align*}
dove \(\hat{S}\) è l'operatore SWAP.
\begin{figure}
	\centering
	\input{img/cns-oCP.quantikz}
	\caption{circuito che mostra la non equivalenza del \emph{conditional non-signalling}
		con la condizione \emph{oCP}.}
	\label{fig:cnsNocp}
\end{figure}
Il circuito è rappresentato in Figura \ref{fig:cnsNocp}.

Il superoperatore da \(r\) a \(t\) è descritto da 
\[\Lambda_{t\colon r} \left[\hat{\rho}^{(r)}\right] =
    \Tr\left[\hat{\rho}^{(r)}\right] \hat{\tau} \]
con \(\hat{\tau} = \frac{1}{2}\Id_\mathcal{S} + \frac{1}{3}\SigmaY_\mathcal{S}\).
Il superoperatore da \(r\) a \(s\) invece è 
\[\Lambda_{s\colon r} \left[\hat{\rho}^{(r)}\right] =
    \Tr\left[\hat{\rho}^{(r)}\right] \frac{\Id_\mathcal{S}}{2}. \]
    
Consideriamo ora il caso in cui interrompiamo l'evoluzione all'istante \(s_-\), e poi
prepariamo il sistema \(\mathcal{S}\) nello stato \(\hat{\rho}^{(s)}\). Nel momento
dell'interruzione l'ambiente risulta essere nello stato
\[\hat{\eta}_{\mathcal{E}_s\mathcal{E}_r} =
      \frac{\Id_{\mathcal{E}_s}\otimes\hat{\rho}^{(r)}_{\mathcal{E}_r}}{2}.\]
Se fissiamo \(\bm{v} = (a,b,c)\) un vettore della Sfera di Bloch, e prepariamo lo
stato 
\[\hat{\rho}^{(s)} =
     \frac{1}{2} \left( \Id_\mathcal{S} + \bm{\sigma}_\mathcal{S} \cdot \bm{v} \right),\]
allora il superoperatore da \(s\) a \(t\) è 
\[\Lambda_{t\colon s} \left[\hat{\rho}^{(s)}\right] =
     \frac{1}{2} \left( \Id_\mathcal{S} + \bm{\sigma}_\mathcal{S} \cdot \bm{w} \right)
     \quad \text{con} \quad \bm{w} = \frac{1}{3} (a-2c, b, 2a-c). \]
 Il processo è chiaramente conditional non-signalling, dato che non c'è informazione
 di \(\hat{\rho}^{(r)}\) nella mappa
\(\Lambda_{s\colon r}\), ma 
\[\Lambda_{t\colon r} \not= \Lambda_{t\colon s} \circ \Lambda_{s\colon r}.\]
\begin{tcolorbox}[colback=red!5!white,colframe=red!75!black,title=Attenzione!]
	Nel paper \cite{CPdoesnotimply} si dice che 
	\(\bm{w} = (\frac{a-2b}{3}, b, \frac{2a-c}{3})\). Secondo me è sbagliato: 
	ad esempio se \(\bm{v} = (0,1,0)\) allora \(\bm{w}\) non sarebbe nella
	sfera di Bloch.
	
	Quello che ho scritto è quanto viene a me, ho fatto verificato
	il conto anche con Python.
\end{tcolorbox}

\subsection[iCP non è equivalente a oCP]
           {iCP \(\centernot=\) oCP} \label{cont:icpocp}
Vedi direttamente \cite{CPdoesnotimply}.

\subsection[oCP non è equivalente a Markoviano]
           {oCP \(\centernot=\) Markov} \label{cont:ocpmar}
Vedi direttamente \cite{CPdoesnotimply}. Per una rappresentazione minimale del
circuito che realizza il controesempio si veda la Figura~\ref{fig:oCPMarkov}.
\begin{figure}
	\centering
	\resizebox{0.6\textwidth}{!}{\input{img/ocp-no-markov.quantikz}}
	\caption{controesmpio oCP diverso da Markov.}
	\label{fig:oCPMarkov}
\end{figure}

\section{Domande senza risposta}
Qui metto tutte le domande che mi sono venute leggendo gli articoli, ma la cui risposta
non è contenuta in essi. Magari semplicemente perché la risposta è talmente
banale che non serviva scriverla.
\paragraph{iCP + cond. non-sign.\(\implies\)oCP?} 
Non riesco proprio a darmi una risposta...

\paragraph{Come si dimostra l'equivalenza tra le due definizioni di markovianità quantistica?}
Non sembra facilissima a prima vista. Basta l'equivalenza classica? Secondo me no, però
insomma, si potrebbe glissare... 

\printbibliography

% Ecco un bell'orso
\clearpage
\section{Extra: un orso}
Non sarebbe una bella presentazione se non ci fosse almeno un easter-egg. Ho scelto
di mettere un'orso abbigliato adeguatamente per l'occasione.
~
\begin{figure}[h!]
	\centering
	\resizebox{0.9\textwidth}{!}{\begin{tikzpicture}
  \bear[graduate=black,tassel=blue]
  \bearwear[shirt=white,
		    body deco={\node at (beartummy)
		    	         {\includegraphics[width=13.5pt]{img/unipi-noscritta.png}};}
           ]
\end{tikzpicture}}
\end{figure}

\end{document}